% 第 15 讲：消没定理蕴含上同调类域论
% 源：src/lecture15.txt（188 行，含 2 处页眉）
\section*{第 15 讲\quad 消没定理蕴含上同调类域论}

上一讲我们把问题重新表述为：对 Galois 群为 $G$ 的非阿基米德局部域扩张 $L/K$，证明
\begin{equation}
(L^{\times})_{tG}\simeq\Z_{tG}[-2]. \tag{15.1}
\end{equation}
于是我们的新目标是计算 $L^{\times}$ 的 Tate 上同调。回忆我们令 $K^{\mathrm{unr}}$ 表示 $K$ 的极大非分歧扩张的完备化；我们希望用 $K^{\mathrm{unr}}$ 来计算这个 Tate 上同调。

\textbf{断言 15.1.} 若 $x\in K^{\mathrm{unr}}$ 在 $K$ 上代数（由于完备化，未必如此），则 $K':=K(x)$ 在 $K$ 上非分歧。

\textit{证明.} 作为 $K$ 的有限代数扩张，$K'$ 是局部域，且有嵌入
\[
\OK_{K'}/\mathfrak{p}_{K}\OK_{K'}\hookrightarrow\OK_{K^{\mathrm{unr}}}/\mathfrak{p}_{K}\OK_{K^{\mathrm{unr}}}=\bar k,
\]
其中 $k:=\OK_{K}/\mathfrak{p}_{K}$。故 $\OK_{K'}/\mathfrak{p}_{K}\OK_{K'}$ 是域，从而 $K$ 与 $K'$ 的素元相同。\hfill$\square$

\textbf{断言 15.2.} $(K^{\mathrm{unr}})^{\sigma=1}=K$，即被 Frobenius 自同构 $\sigma\in G$ 固定（也就是在其作用下为恒等）的元素恰为 $K$；$\sigma$ 由各个非分歧扩张的 Frobenius 经连续性延拓到完备化而得到。

回忆我们有短正合列
\[
0\to K\to K^{\mathrm{unr}}\xrightarrow{\ 1-\sigma\ }K^{\mathrm{unr}},
\]
它在乘法群上可改写为
\[
1\to K^{\times}\to K^{\mathrm{unr},\times}\xrightarrow{\ x\mapsto x/\sigma x\ }K^{\mathrm{unr},\times}\xrightarrow{\ v\ }\Z\to 0.
\]
我们已证明：$K^{\mathrm{unr},\times}$ 中的元素若能写成 $x/\sigma x$，则它必是整数环中的单位 $x\in\OK_{K^{\mathrm{unr}}}^{\times}$；该映射在每个伴随分次项上是同构，从而在 $\OK_{K^{\mathrm{unr}}}^{\times}$ 上是同构。

现在，我们想显式构造 (15.1) 中的同构。第一次尝试如下：设 $L/K$ 完全分歧（既然上一讲已处理非分歧情形，这个假设相当温和），于是 $L^{\mathrm{unr}}=L\otimes_{K}K^{\mathrm{unr}}$。则我们有如下定理（稍后证明）。

\textbf{定理 15.3（消没定理）.} 若 $L/K$ 完全分歧，则复形 $(L^{\mathrm{unr},\times})_{tG}$ 无环。

\textbf{断言 15.4.} 消没定理蕴含上同调类域论。

\textit{证明.} 设 $L/K$ 完全分歧。我们有四项正合列
\begin{equation}
1\to L^{\times}\to L^{\mathrm{unr},\times}\xrightarrow{\ x\mapsto x/\sigma x\ }L^{\mathrm{unr},\times}\xrightarrow{\ v\ }\Z\to 0. \tag{15.2}
\end{equation}
可将其改写如下：
\[
\begin{array}{c|cccccc}
A & \cdots & 0 & L^{\times} & 0 & 0 & \cdots\\[2pt]
B & \cdots & 0 & L^{\mathrm{unr},\times} & L^{\mathrm{unr},\times} & 0 & \cdots\\[2pt]
\operatorname{Coker}(A\to B) & \cdots & 0 & 0 & \Z & 0 & \cdots
\end{array}
\]
其中 $L^{\times}$ 在 $-1$ 次，$B$ 的微分是 $1-\sigma$、$\operatorname{Coker}(A\to B)$ 的微分是 $v$。最后到由 (15.2) 所得同伦余核的拟同构来自断言 10.12，因为 $A\hookrightarrow B$ 是单射（注意这对任何四项正合列都成立）。逐项余核给出单射
\[
L^{\mathrm{unr},\times}/L^{\times}\xrightarrow{\ x\mapsto x/\sigma x\ }L^{\mathrm{unr},\times},
\]
因为去掉商之后，$L^{\times}$ 恰是该映射的核。

现在，我们有拟同构
\[
B_{tG}=\operatorname{hCoker}\bigl(L^{\mathrm{unr},\times}\xrightarrow{\ 1-\sigma\ }L^{\mathrm{unr},\times}\bigr)_{tG}\simeq\operatorname{hCoker}\bigl((L^{\mathrm{unr},\times})_{tG}\to(L^{\mathrm{unr},\times})_{tG}\bigr),
\]
于是由消没定理，$(L^{\mathrm{unr},\times})_{tG}$ 无环，从而由上同调长正合列，这个同伦余核也无环。因此
\[
(L^{\times}[2])_{tG}=\operatorname{hCoker}\bigl((L^{\times}[1])_{tG}\to 0\bigr)=\operatorname{hCoker}\bigl(A_{tG}\to B_{tG}\bigr)\simeq\Z_{tG},
\]
正是所要。\hfill$\square$

现在设 $L/K$ 是一般的有限 Galois 扩张，$G:=\Gal(L/K)$（虽然我们可以分别处理非分歧与完全分歧情形，因为任何扩张都典范地是这样的扩张的复合）。若 $L/K$ 非分歧，则
\[
L\otimes_{K}K^{\mathrm{unr}}=\prod_{L\hookrightarrow K^{\mathrm{unr}}}K^{\mathrm{unr}}
\]
典范地，按这些嵌入编号。事实上，下述结论成立：

\textbf{定理 15.5（一般消没定理）.} $\bigl[(L\otimes_{K}K^{\mathrm{unr}})^{\times}\bigr]_{tG}$ 无环。

为理解 $L\otimes_{K}K^{\mathrm{unr}}$ 的结构，注意 $\widehat{\Z}\sigma$ 作用在第二个因子、$G$ 作用在第一个因子；这两个作用（即 $x\otimes y\mapsto gx\otimes y$ 与 $x\otimes y\mapsto x\otimes\sigma y$）显然交换。同样，被 $\sigma$ 固定的点构成
\[
L=L\otimes_{K}K\hookrightarrow L\otimes_{K}K^{\mathrm{unr}}.
\]

\textbf{断言 15.6.} 下述序列正合：
\[
1\to L^{\times}\to(L\otimes_{K}K^{\mathrm{unr}})^{\times}\xrightarrow{\ x\mapsto x/\sigma x\ }(L\otimes_{K}K^{\mathrm{unr}})^{\times}\to\Z\to 0.
\]

\textit{证明.} 若 $x\in K^{\mathrm{unr}}$ 是单位，则 $\sigma x$ 也是，故映射 $x\mapsto x/\sigma x$ 良定义；并且 $x$ 在其核中当且仅当 $x$ 在 $\sigma$ 作用下不动，即 $x\in K$；又 $L\otimes_{K}K=L$，故得到 $L$ 的一个单位，这证明了左半部分的正合性。至于到 $\Z$ 的映射，由
\[
\begin{CD}
(L\otimes_{K}K^{\mathrm{unr}})^{\times} @>>> \Z\\
@V{N_{L/K}\otimes\operatorname{id}}VV @A{\ v\ }AA\\
K\otimes_{K}K^{\mathrm{unr},\times} @= K^{\mathrm{unr},\times}
\end{CD}
\]
定义，其中 $N_{L/K}(x):=\prod_{g\in G}gx$。于是其核为 $\OK_{K^{\mathrm{unr}}}^{\times}$，它恰是 $x\mapsto x/\sigma x$ 的像。此外该映射是满射，因为 $1\otimes\pi\mapsto 1$。\hfill$\square$

注意若 $L/K$ 完全分歧，这就回到我们先前的扩张。事实上，若记 $L^{\mathrm{unr}}=L\otimes_{K}K^{\mathrm{unr}}$，则两边的 $\sigma$「吻合」，即 $L^{\mathrm{unr}}$ 与 $K^{\mathrm{unr}}$ 上的诱导 Frobenius 自同构相同，因为 $L$ 与 $K$ 有相同的剩余域。该扩张的范映射 $N_{L/K}\colon L^{\mathrm{unr},\times}\to K^{\mathrm{unr},\times}$ 满足 $v_{K^{\mathrm{unr}}}\circ N=v_{L^{\mathrm{unr}}}$（这样的扩张由 $K$ 的素元的 $n$ 次根生成，于是 $N(\pi^{1/n})=\pi$）。

现在设 $L/K$ 是 $n$ 次非分歧扩张。固定一个嵌入 $L\hookrightarrow K^{\mathrm{unr}}$，并用 $\sigma\in\Gal(L/K)$ 同时表示 $L/K$ 的 Frobenius 元素。则我们有同构
\[
L\otimes_{K}K^{\mathrm{unr}}\xrightarrow{\ \sim\ }\prod_{i=0}^{n-1}K^{\mathrm{unr}},\qquad x\otimes y\mapsto\bigl((\sigma^{i}x)\cdot y\bigr)_{i=0}^{n-1},
\]
其中乘积经由我们固定的嵌入（注意也可以像先前那样对嵌入取乘积，从而更典范地做到这一点）。现在问：$L\otimes_{K}K^{\mathrm{unr}}$ 的自同构 $\mathrm{id}\otimes\sigma$ 在该同构下对应什么？我们有
\[
x\otimes\sigma y\mapsto(x\cdot\sigma y,\ \sigma x\cdot\sigma y,\ \sigma^{2}x\cdot\sigma y,\ \dots)=\sigma\bigl(\sigma^{-1}x\cdot y,\ x\cdot y,\ \sigma x\cdot y,\ \dots\bigr),
\]
所以它是 $\sigma$ 作用在 $x\otimes y$ 的像的「右旋」上（注意 $\sigma$ 在 $K^{\mathrm{unr}}$ 上不是有限阶的，$\mathrm{id}\otimes\sigma$ 也不应是有限阶的，这就排除了把 $\mathrm{id}\otimes\sigma$ 的像写成旋转的可能性）。类似地，范映射 $N_{L/K}\colon\prod K^{\mathrm{unr},\times}\to K^{\mathrm{unr},\times}$ 取所有分量的乘积。

我们希望某个元素 $(x_{0},\dots,x_{n-1})\in\prod K^{\mathrm{unr},\times}$ 落在 $y/\sigma y$ 的像中（即断言 15.6 正合列中间那个映射；这里 $\sigma$ 指自同构 $\mathrm{id}\otimes\sigma$），当且仅当 $\prod x_{i}\in\OK_{K^{\mathrm{unr}}}^{\times}$，即 $\sum v(x_{i})=0$。回忆反向蕴含是平凡的，因为我们已经证明 $\OK_{K^{\mathrm{unr}}}^{\times}\xrightarrow{\ y/\sigma y\ }\OK_{K^{\mathrm{unr}}}^{\times}$ 是满射，因为它在伴随分次层面满射。至于正向，我们有
\[
(y_{0},\dots,y_{n-1})\xrightarrow{\ y/\sigma y\ }\Bigl(\frac{y_{0}}{\sigma y_{n-1}},\frac{y_{1}}{\sigma y_{0}},\dots\Bigr)=:(x_{0},x_{1},\dots).
\]
于是
\[
y_{0}=x_{0}\cdot\sigma y_{n-1},
\]
\[
y_{1}=x_{1}\cdot\sigma y_{0}=x_{1}\cdot\sigma x_{0}\cdots\sigma^{2}y_{n-1},
\]
\[
\cdots=\cdots
\]
\[
y_{n-1}=x_{n-1}\cdot\sigma x_{n-2}\cdots\sigma^{n-1}x_{0}\cdot\sigma^{n}y_{n-1},
\]
即
\[
\frac{y_{n-1}}{\sigma^{n}y_{n-1}}=x_{n-1}\cdot\sigma x_{n-2}\cdots\sigma^{n-1}x_{0}.
\]
注意这里一切都是 $K^{\mathrm{unr}}$ 的元素，所以我们确实没有 $\sigma^{n}=\mathrm{id}$！上一讲我们证明了这当且仅当右端属于 $\OK_{K^{\mathrm{unr}}}^{\times}$，而这等价于说 $\sum v(x_{i})=0$。该正合列的一般情形是这两种情形的某种混合。

我们现在把这些结果与上一讲的结果作比较。假定消没定理成立。对非分歧扩张 $L/K$，我们有两个从 $(L^{\times})_{tG}$ 到 $\Z[-2]_{tG}$ 的拟同构：一个来自我们刚刚所做，另一个来自 $(\OK_{L}^{\times})_{tG}\simeq 0$ 蕴含（由循环性）$(L^{\times})_{tG}\simeq\Z_{tG}\simeq(\Z[-2])_{tG}$。我们断言这两个拟同构重合。证明的思路如下：$G=\Z/n\Z$（生成元为 Frobenius 元素），且有短正合列
\[
0\to\Z\to\Z[G]\xrightarrow{\ 1-\sigma\ }\Z[G]\to\Z\to 0.
\]
如习题集 7 第 1(e) 题所示，$\Z[G]_{tG}\simeq 0$ 是拟同构（这很容易证明，而且我们已经对循环群证明过）。于是得到 $\Z_{tG}[2]\simeq\Z_{tG}$，我们断言这与由复形的 $2$-周期性得到的同构相同。证明靠图表追踪。我们有 $(L\otimes_{K}K^{\mathrm{unr}})^{\times}=\prod K^{\mathrm{unr},\times}$，是有限乘积。于是图表
\[
1\longrightarrow L^{\times}\longrightarrow(L\otimes_{K}K^{\mathrm{unr}})^{\times}\xrightarrow{\ x\mapsto x/\sigma x\ }(L\otimes_{K}K^{\mathrm{unr}})^{\times}\xrightarrow{\ v\ }\Z\longrightarrow 0
\]
与
\[
1\longrightarrow\Z\longrightarrow\Z[G]\xrightarrow{\ 1-\sigma\ }\Z[G]\xrightarrow{\ \varepsilon\ }\Z\longrightarrow 0
\]
竖直映射依次取 $v$、$\prod v$、$\sum v$ 时交换，其中 $\varepsilon$ 表示 $\Z[G]$ 各坐标之和。这恰好说明由两个四项正合列得到的同构相同。

要点是：在 lcft 之下，我们有同构
\[
K^{\times}/N(L^{\times})\cong\Z/n\Z,\qquad\pi\mapsto\mathrm{Frob}.
\]
于是我们把 lcft 归结为消没定理，后者我们将在下一讲证明。
